\documentclass[11pt]{article}
\usepackage[margin=1in]{geometry}
\usepackage{times}
\usepackage{natbib}
\usepackage{graphicx}
\usepackage{booktabs}
\usepackage{hyperref}
\usepackage{amsmath}
\usepackage{xcolor}
\usepackage{longtable}

\hypersetup{colorlinks=true, linkcolor=blue, citecolor=blue, urlcolor=blue}

\newcommand{\pending}[1]{\textcolor{red!70!black}{\textbf{[PENDING: #1]}}}

\title{Citation-Grounded Classification from Raw Network Flow Features:\\
Claude Sonnet versus GPT-4o-mini on the Real CIC-IoT-2023 Dataset}
\author{Gangadhar Singh Shiva \and Poonam Singh Umakanth}
\date{\today}

\begin{document}
\maketitle

\begin{abstract}
Citation-grounded classification pipelines reduce hallucination by requiring
a model to justify a security finding with a reference to an authoritative
external identifier (a CAPEC attack pattern or CVE record) rather than an
unsupported assertion. Prior work in this line measured citation accuracy
when the model is given a natural-language description of an observed
finding, either hand-written or derived from an image. This paper removes
that natural-language intermediary entirely: each trial presents a model
with the raw, 39-feature numeric flow vector extracted directly from a real
network capture, with no interpretive paraphrase, and asks it to classify
and cite. We sample 1,360 flows (40 per class, fixed seed) from the real
CIC-IoT-2023 dataset (34 classes, 46{,}776{,}697 rows, 8.4~GB, obtained and
verified locally) and compare Claude Sonnet and GPT-4o-mini under two
context-construction strategies (full-corpus and TF-IDF top-5 retrieval)
against the same 15-entry, individually verified CAPEC/CVE corpus used
throughout this project's prior work. Nineteen of the 34 real classes have
a directly matching real corpus entry (ground-truth trials); the remaining
15, including benign traffic, do not and are legitimate
abstention-expected trials. Of 5,440 planned real API calls, 387 (all
Claude, 14.2\% of Claude's total) failed with a disclosed, non-modeling
cause -- an Anthropic account usage-limit error -- and are excluded from
every metric in this report rather than counted as errors; critically,
these failures are \emph{not} randomly distributed across classes (
Section~\ref{sec:limitations}) and this is reported rather than smoothed
over. On the 5,053 successfully scored real trials: correct-citation rate
was 52.1\%~[48.1,56.0] (Claude, full-corpus) and 50.2\%~[46.2,54.1] (Claude,
RAG) versus 36.7\%~[33.3,40.3] (GPT-4o-mini, full-corpus) and
36.8\%~[33.4,40.3] (GPT-4o-mini, RAG) -- a large, statistically significant
gap favoring Claude at fixed strategy (McNemar $p<0.0001$ full-corpus,
$p=0.0007$ RAG). Zero invented IDs occurred across all 5,053 trials. Correct
abstention on the 15 no-match classes was rare for both models
(0.3\%--2.1\%), a substantially worse abstention profile than this
project's earlier natural-language-query experiments on the same corpus.
Per-class results show a sharp split: flood-type volumetric attacks
(TCP/UDP/ICMP/SYN floods) were classified with 87.5--100\% accuracy by both
models, while protocol-semantic and application-layer attacks (SQL
injection, XSS, brute-force, Mirai infection, spoofing) scored at or near
0\% -- evidence that a 39-feature flow summary alone carries strong signal
for volumetric attacks but effectively none for attacks whose signature is
in payload content or multi-flow behavioral context. We report this as a
disclosed partial dataset, not a completed one: 387 Claude trials remain
pending an account usage-limit reset and are explicitly out of scope for
this report's numbers.
\end{abstract}

\section{Introduction}

Citation-grounded architectures for automated threat classification reduce
one specific failure mode -- a model asserting a finding with no traceable
external evidence -- by requiring every classification to reference a real,
verifiable identifier drawn from a fixed corpus, or to explicitly decline
when no entry applies. This project's prior work measured that architecture
under two conditions: a purely synthetic, hand-paraphrased query set
(\citet{clearsight2026}, the full-corpus-vs-RAG retrieval study), and a
natural-language "observed finding" description manually or programmatically
derived from an underlying artifact (the Claude-vs-GPT-4o-mini
citation-accuracy study). In both cases, a human- or pipeline-authored
English sentence stood between the raw evidence and the model.

That intermediary is a confound. A natural-language description of a
network flow already encodes an interpretation -- "sustained high-rate SYN
packets with no completed handshake" is itself a classification, expressed
in different words. A model shown that sentence is not being asked to
classify raw evidence; it is being asked to recognize a paraphrase of the
answer. This paper removes the paraphrase. Each trial in this experiment
shows the model the literal 39-feature numeric vector extracted from one
real network flow in the CIC-IoT-2023 dataset~\citep{neto2023ciciot2023} --
header length, protocol type, TTL, per-flag packet counts, protocol
indicator flags, and packet-size/inter-arrival-time statistics -- with no
interpretive text of any kind, and asks the model to classify the flow and
cite the matching entry in a fixed, real CAPEC/CVE grounding corpus, or to
abstain if none applies.

This is a harder and more realistic task than the natural-language version:
it is what a citation-grounded pipeline would actually face if it ingested
NetFlow-style records directly, without a hand-authored intermediate
summary (cf.\ multimodal approaches that instead visualize flow features as
images before classification, e.g.\ \citet{netflowviz}). It also lets this
paper measure something the earlier natural-language experiments could
not: whether citation accuracy on this project's corpus is uniform across
attack types, or concentrated in the attack types whose statistical
signature survives being reduced to 39 numbers.

Contributions:
\begin{enumerate}
    \item A real-data sampling protocol (Section~\ref{sec:dataset}) drawn
    directly from the full, locally verified CIC-IoT-2023 CSV release (34
    classes, 46.7 million rows, 8.4~GB) -- not a synthetic or
    previously-subsampled derivative.
    \item A disclosed, programmatically derived ground-truth mapping
    (Section~\ref{sec:corpus}) from 19 of the 34 real classes to this
    project's existing 15-entry verified CAPEC/CVE corpus, and an explicit
    accounting of the 15 classes with no real matching entry.
    \item A raw-feature task design (Section~\ref{sec:task}) that shows the
    model literal extracted features rather than a natural-language
    paraphrase, isolating citation accuracy from paraphrase-recognition.
    \item Real results (Section~\ref{sec:results}) from 5,053 successfully
    scored live API trials across two models and two context-construction
    strategies, including a per-class breakdown that surfaces a large,
    previously unmeasured accuracy split between volumetric and
    protocol-semantic attack types.
    \item A precise, non-aggregated account of a real data-collection
    failure (Section~\ref{sec:limitations}): 387 Claude trials are missing
    because of an account usage limit, and -- unlike a simple "14.2\% of
    trials failed" disclosure -- we show these failures are concentrated in
    the last five classes processed, so four real classes currently have
    \emph{zero} Claude data, a materially different situation from a
    uniform random 14.2\% loss.
\end{enumerate}

\section{Related Work}

\citet{clearsight2026} introduced the citation-grounded architecture this
paper's task design builds on. The companion full-corpus-vs-RAG retrieval
study measured retrieval recall, cost, and latency across TF-IDF,
BM25~\citep{robertson2009bm25}, and LSA~\citep{deerwester1990lsa} backends
against a synthetic query set, without exercising a live classifying model.
A follow-up study exercised live Claude and GPT-4o-mini API calls on a
45-query, hand-paraphrased natural-language version of the same task,
finding zero invented IDs across 180 scored trials and a
context-strategy-sensitivity asymmetry between the two models. This paper
is the natural-data extension of that line: the same corpus and the same
two models, but the query is now a real, unmodified flow-feature vector
rather than an authored sentence.

Multimodal approaches to the same underlying problem visualize NetFlow
features as images and classify the image directly
\citep{netflowviz}, trading the raw-feature representation used here for a
visual one; this project's separate CAMLIS-track paper compares that visual
ingestion pathway against text-native ingestion for the same dataset. The
CIC-IoT-2023 dataset itself~\citep{neto2023ciciot2023} provides labeled
traffic from 33 attack types across 7 categories plus benign traffic,
captured from a 105-device IoT topology; this paper uses all 34 provided
labels. Anomaly detection surveys~\citep{chandola2009anomaly} and classical
statistical/ML detectors form the pre-LLM baseline this citation-grounded
line of work is positioned against: those detectors flag a deviation but do
not produce a human-auditable, externally verifiable citation for it, which
is the specific gap this architecture targets.

\section{Dataset}
\label{sec:dataset}

\subsection{Source}

All results in this paper use the CIC-IoT-2023
dataset~\citep{neto2023ciciot2023}, CSV release, obtained and verified
directly (not a third-party mirror or resample). The local copy contains
309 CSV files across 34 class subdirectories, totaling 8.4~GB and
46{,}776{,}697 data rows (row counts confirmed by direct line-count over
every file, not read from documentation). Each row is one extracted network
flow with 39 numeric features: \texttt{Header\_Length}, \texttt{Protocol
Type}, \texttt{Time\_To\_Live}, \texttt{Rate}, seven per-flag indicator/
count columns (SYN/FIN/RST/PSH/ACK/ECE/CWR and their counts), eleven
protocol-presence indicator columns (HTTP, HTTPS, DNS, Telnet, SMTP, SSH,
IRC, TCP, UDP, DHCP, ARP, ICMP, IGMP, IPv, LLC), and eight packet-size/
timing statistics (total sum, min, max, average, standard deviation, total
size, inter-arrival time, packet count, variance).

Table~\ref{tab:classcounts} shows the real per-class row-count range: from
1{,}252 rows (\texttt{Uploading\_Attack}, the smallest class) to
7{,}200{,}501 rows (\texttt{DDoS-ICMP\_Flood}, the largest) -- a
5{,}750$\times$ imbalance across classes. Because even the smallest class
has over a thousand rows, this paper's sampling design (below) is not
data-constrained the way an earlier, image-based experiment on the same
project's visualization dataset was (that dataset had only 4--5 images per
class).

\begin{table}[h]
\centering
\small
\begin{tabular}{lr}
\toprule
\textbf{Statistic} & \textbf{Value} \\
\midrule
Classes & 34 (33 attack + Benign) \\
CSV files & 309 \\
Total size & 8.4~GB \\
Total data rows & 46{,}776{,}697 \\
Smallest class (rows) & 1{,}252 (Uploading\_Attack) \\
Largest class (rows) & 7{,}200{,}501 (DDoS-ICMP\_Flood) \\
Features per row & 39 \\
\bottomrule
\end{tabular}
\caption{Real, directly measured statistics of the local CIC-IoT-2023 CSV
dataset used in this paper.}
\label{tab:classcounts}
\end{table}

\subsection{Sampling}

We draw a stratified sample of $N{=}40$ rows per class (1{,}360 rows total),
fixed seed 42, using memory-safe reservoir-style index sampling: for each
class, global row indices are drawn uniformly without replacement across
that class's constituent files (weighted correctly by each file's row
count), then each source file is scanned once to extract only the sampled
rows -- no file is ever fully loaded into memory. $N{=}40$/class was chosen
over a larger $N$ (e.g.\ 100/class) as a cost/precision balance: at
$N{=}40$, a Wilson confidence interval on an $\approx$85\% accuracy
estimate is roughly $\pm$11 points, versus $\pm$16 points at $N{=}20$, for
an incremental real-API cost of approximately \$10 -- judged worthwhile;
$N{=}100$ would tighten this further but at roughly 2.5$\times$ the cost for
diminishing precision gain. The resulting manifest (sample id, true class,
ground-truth citation id, source file, row index, and the exact serialized
query text -- see Section~\ref{sec:task}) is a frozen, versioned artifact:
the experiment-runner script never re-touches the source CSVs, so the exact
sampled rows are reproducible and auditable independent of the multi-GB
source files.

\section{Grounding Corpus and Ground-Truth Mapping}
\label{sec:corpus}

This paper reuses, unchanged, the 15-entry real CAPEC/CVE grounding corpus
established and individually verified (each ID fetched from its official
source, \url{capec.mitre.org} or \url{nvd.nist.gov}) in this project's
prior retrieval and citation-accuracy studies: 12 CAPEC attack patterns and
3 CVE records (Table~\ref{tab:corpus}).

\begin{table}[h]
\centering
\small
\begin{tabular}{lll}
\toprule
\textbf{ID} & \textbf{Name} & \textbf{Kind} \\
\midrule
CAPEC-482 & TCP Flood & CAPEC \\
CAPEC-488 & HTTP Flood & CAPEC \\
CAPEC-486 & UDP Flood & CAPEC \\
CAPEC-487 & ICMP Flood & CAPEC \\
CAPEC-300 & Port Scanning & CAPEC \\
CAPEC-49 & Password Brute Forcing & CAPEC \\
CAPEC-94 & Adversary in the Middle (AiTM) & CAPEC \\
CAPEC-141 & Cache Poisoning & CAPEC \\
CAPEC-66 & SQL Injection & CAPEC \\
CAPEC-63 & Cross-Site Scripting (XSS) & CAPEC \\
CAPEC-137 & Parameter Injection & CAPEC \\
CAPEC-153 & Input Data Manipulation & CAPEC \\
CVE-2016-10401 & ZyXEL PK5001Z default \texttt{su} password & CVE \\
CVE-2017-17215 & Huawei HG532 remote code execution & CVE \\
CVE-2014-9583 & ASUS WRT \texttt{infosvr} backdoor & CVE \\
\bottomrule
\end{tabular}
\caption{Grounding corpus, unchanged from prior work in this project.}
\label{tab:corpus}
\end{table}

Each corpus entry carries a list of the real CIC-IoT-2023 class names it
covers (its \texttt{attack\_families} field, established when this corpus
was first built). We use that existing, disclosed field -- not a mapping
hand-picked for this paper -- to programmatically determine, for each of
the 34 real classes, whether a real ground-truth citation id exists.
Table~\ref{tab:mapping} shows the result: 19 of 34 classes map to a real
entry; the remaining 15, including \texttt{Benign\_Final}, have no matching
entry and are legitimate abstention-expected trials, exactly analogous to
the held-out abstention queries in this project's earlier natural-language
experiment. This is not a gap papered over with synthetic entries: it is
the honest coverage of a small, hand-verified real corpus, and the 15
unmapped classes are scored on whether the model correctly recognizes that
nothing in the corpus applies, not silently dropped from the experiment.

\begin{table}[h]
\centering
\small
\begin{tabular}{ll}
\toprule
\textbf{19 classes with a real ground-truth id} & \textbf{Mapped entry} \\
\midrule
DDoS-HTTP\_Flood, DoS-HTTP\_Flood & CAPEC-488 \\
DDoS-ICMP\_Flood & CAPEC-487 \\
DDoS-SYN\_Flood, DDoS-TCP\_Flood, DoS-SYN\_Flood, DoS-TCP\_Flood & CAPEC-482 \\
DDoS-UDP\_Flood, DoS-UDP\_Flood & CAPEC-486 \\
DNS\_Spoofing & CAPEC-141 \\
DictionaryBruteForce & CAPEC-49 \\
MITM-ArpSpoofing & CAPEC-94 \\
Mirai-greeth\_flood, Mirai-greip\_flood & CVE-2017-17215 \\
Mirai-udpplain & CVE-2014-9583 \\
Recon-PortScan & CAPEC-300 \\
SqlInjection & CAPEC-66 \\
Uploading\_Attack & CAPEC-137 \\
XSS & CAPEC-63 \\
\midrule
\textbf{15 classes with no real match (abstention-expected)} & \\
\midrule
\multicolumn{2}{l}{Backdoor\_Malware, Benign\_Final, BrowserHijacking, CommandInjection,} \\
\multicolumn{2}{l}{DDoS-ACK\_Fragmentation, DDoS-ICMP\_Fragmentation, DDoS-PSHACK\_FLOOD,} \\
\multicolumn{2}{l}{DDoS-RSTFINFLOOD, DDoS-SlowLoris, DDoS-SynonymousIP\_Flood,} \\
\multicolumn{2}{l}{DDoS-UDP\_Fragmentation, Recon-HostDiscovery, Recon-OSScan,} \\
\multicolumn{2}{l}{Recon-PingSweep, VulnerabilityScan} \\
\bottomrule
\end{tabular}
\caption{Programmatic ground-truth mapping from real CIC-IoT-2023 classes
to the 15-entry corpus, derived from each entry's existing
\texttt{attack\_families} field. \texttt{Mirai-udpplain} is also plausibly
covered by CVE-2016-10401 (ZyXEL); the mapping resolves such ties
deterministically to one canonical id, a disclosed simplification (multiple
real CVEs are plausible root causes for the same observed traffic class).}
\label{tab:mapping}
\end{table}

\section{Task Design}
\label{sec:task}

\subsection{Raw-feature query construction}

Each trial's query is a single, literal serialization of one sampled row's
39 feature values as \texttt{name=value} pairs (e.g.\ \texttt{Header\_
Length=6.4, Protocol Type=17, Time\_To\_Live=60.2, Rate=35.06, ...}),
prefixed with the fixed string ``Observed network flow, extracted feature
vector:''. No natural-language interpretation, protocol name lookup (e.g.\
mapping the numeric protocol id 17 to ``UDP''), or statistical summary is
added. This is a deliberate departure from this project's earlier
natural-language query design: it removes the possibility that the model is
merely pattern-matching a paraphrase, at the cost of a harder, more
realistic task, and it is the central methodological difference this paper
studies relative to that earlier work.

\subsection{Response format and grading}

Unchanged from this project's established citation-grounded task
definition: the model is shown a set of candidate reference entries (either
all 15 corpus entries, or the top-5 TF-IDF-retrieved entries, depending on
condition) and the single query, and must respond with (a) a short
classification label, (b) a cited entry id drawn only from the entries
shown or \texttt{null} if none applies, and (c) a one-sentence
justification. Grading, also unchanged: \textbf{correct-citation} (cited id
equals the true class's mapped ground-truth id, ground-truth trials only),
\textbf{invented-id} (cited id was not among the entries actually shown),
and \textbf{correctly-abstained} (cited id is null, abstention-expected
trials only).

\subsection{Models and context-construction strategies}

Two models: \textbf{Claude} (Anthropic API, model string
\texttt{claude-sonnet-5}) and \textbf{GPT-4o-mini} (OpenAI API, model
string \texttt{gpt-4o-mini}). Two context-construction strategies, reusing
the \texttt{FullCorpusStrategy} and \texttt{RAGStrategy} (TF-IDF, $k{=}5$)
implementations from this project's prior retrieval study unchanged:
\textbf{full-corpus} (all 15 entries shown every trial) and \textbf{RAG}
(top-5 TF-IDF-retrieved entries only). As in prior work, Claude's API
rejects an explicit \texttt{temperature} parameter for this model string
(a disclosed deviation carried over from the earlier study); GPT-4o-mini
runs at temperature~0. Planned trial count: 1{,}360 samples $\times$ 2
models $\times$ 2 strategies $=$ 5{,}440 real API calls.

\subsection{Resumable execution}

Because a 5{,}440-call real-money run can be interrupted by rate limits,
network faults, or (as happened in this run) an account usage cap, the
harness checkpoints every trial result to disk immediately upon completion
and, on any restart, skips every (sample, model, strategy) combination that
already has a successful result, retrying only what previously failed or
never ran. This is what makes it possible to report a disclosed partial
result in Section~\ref{sec:results} rather than being forced to either wait
indefinitely or discard already-collected, already-paid-for data.

\section{Statistical Plan}
\label{sec:stats}

Point estimates use bootstrap 95\% confidence intervals (10{,}000
resamples), consistent with this project's prior work. Because both models
are run on the same sampled rows, the Claude-vs-GPT-4o-mini comparison at
fixed strategy is paired, and a continuity-corrected McNemar's test is used
on correct-citation outcomes; the same pairing and test are used for
full-corpus-vs-RAG at fixed model.

\section{Results}
\label{sec:results}

\subsection{Scope of this report}

Of 5{,}440 planned real trials, 5{,}053 completed successfully; 387 (all
Claude, 193 full-corpus and 194 RAG) failed with the same disclosed,
non-modeling cause: \texttt{"You have reached your specified API usage
limits. You will regain access on 2026-09-01."} Per this paper's explicit
scope decision, these 387 trials are excluded from every metric below --
not imputed, not counted as incorrect, not waited on. Section~\ref{sec:limitations}
gives the precise, non-uniform distribution of these failures across
classes, which materially affects how the Claude numbers below should be
read.

\subsection{Citation accuracy}

\begin{table}[h]
\centering
\small
\begin{tabular}{llrccc}
\toprule
\textbf{Model} & \textbf{Strategy} & \textbf{N} & \textbf{Correct-citation} & \textbf{Invented-ID} & \textbf{Correct-abstain} \\
\midrule
Claude & full-corpus & 607 & 0.521~[0.481,0.560] & 0.000~[0.000,0.000] & 0.021~[0.011,0.034] \\
Claude & RAG ($k{=}5$) & 606 & 0.502~[0.462,0.541] & 0.000~[0.000,0.000] & 0.016~[0.007,0.027] \\
GPT-4o-mini & full-corpus & 760 & 0.367~[0.333,0.403] & 0.000~[0.000,0.000] & 0.005~[0.000,0.012] \\
GPT-4o-mini & RAG ($k{=}5$) & 760 & 0.368~[0.334,0.403] & 0.000~[0.000,0.000] & 0.003~[0.000,0.008] \\
\bottomrule
\end{tabular}
\caption{Citation accuracy, real live API results, 5{,}053 successfully
scored trials (387 excluded per Section~\ref{sec:limitations}). N is the
ground-truth-trial count actually scored (out of a planned 760 per cell for
the 19 mapped classes); correct-abstain N is separately computed over the
abstention-expected trials (planned 600 per cell) and not shown as a
column here for space, see \texttt{csv\_flow\_summary.csv}. Bootstrap 95\%
CIs, 10{,}000 resamples.}
\label{tab:results-main}
\end{table}

Zero invented IDs across all 5{,}053 scored trials, both models, both
strategies -- consistent with every prior experiment in this project.
Correct-citation rate is substantially lower than the earlier
natural-language-query experiment on the same corpus (which measured
82--96\%): Claude scores 50--52\%, GPT-4o-mini 37\%, on this raw-feature
version of the task. Correct-abstention is rare for both models
(0.3--2.1\%), markedly worse than the earlier natural-language experiment's
abstention rates.

\subsection{Paired comparisons}

\begin{table}[h]
\centering
\small
\begin{tabular}{lccc}
\toprule
\textbf{Strategy} & \textbf{$n$ pairs} & \textbf{$b$, $c$} & \textbf{$\chi^2$ (p-value)} \\
\midrule
full-corpus & 607 & $b{=}44$, $c{=}7$ & $\chi^2{=}25.412$ ($p{<}0.0001$) \\
RAG ($k{=}5$) & 606 & $b{=}35$, $c{=}11$ & $\chi^2{=}11.500$ ($p{=}0.0007$) \\
\bottomrule
\end{tabular}
\caption{McNemar's test, Claude vs.\ GPT-4o-mini, correct-citation, paired
by sample, at fixed strategy.}
\label{tab:mcnemar-model}
\end{table}

\begin{table}[h]
\centering
\small
\begin{tabular}{lccc}
\toprule
\textbf{Model} & \textbf{$n$ pairs} & \textbf{$b$, $c$} & \textbf{$\chi^2$ (p-value)} \\
\midrule
Claude & 606 & $b{=}18$, $c{=}6$ & $\chi^2{=}5.042$ ($p{=}0.0247$) \\
GPT-4o-mini & 760 & $b{=}1$, $c{=}2$ & $\chi^2{=}0.000$ ($p{=}1.0000$) \\
\bottomrule
\end{tabular}
\caption{McNemar's test, full-corpus vs.\ RAG, correct-citation, paired by
sample, at fixed model.}
\label{tab:mcnemar-strategy}
\end{table}

Claude's correct-citation rate is significantly higher than GPT-4o-mini's
at both strategies ($p{<}0.0001$ full-corpus, $p{=}0.0007$ RAG) -- unlike
the earlier natural-language experiment, where the two models were
statistically indistinguishable. Within-model, full-corpus outperforms RAG
significantly for Claude ($p{=}0.0247$) but not for GPT-4o-mini
($p{=}1.0000$, $b$ and $c$ are both near zero, i.e.\ almost no samples flip
outcome between strategies).

\subsection{Per-class breakdown}

Table~\ref{tab:byclass} groups the 19 ground-truth classes by observed
accuracy pattern (full detail in \texttt{csv\_flow\_by\_class.csv}). The
split is stark: every flood-type volumetric class (SYN/TCP/UDP/ICMP floods,
both DoS and DDoS variants) scores 87.5--100\% for both models, while every
protocol-semantic or application-layer class (SQL injection, XSS, brute
force, port scanning, ARP/DNS spoofing, all three Mirai infection vectors)
scores at or near 0\% for both models. \texttt{DDoS/DoS-HTTP\_Flood} is the
one mapped class that falls in between (0--55\%, the most model- and
strategy-sensitive class in the corpus).

\begin{table}[h]
\centering
\small
\begin{tabular}{llc}
\toprule
\textbf{Pattern} & \textbf{Classes} & \textbf{Accuracy range (both models)} \\
\midrule
Volumetric flood & DDoS/DoS-\{SYN,TCP,UDP,ICMP\}\_Flood & 0.875--1.000 \\
HTTP flood (mixed) & DDoS/DoS-HTTP\_Flood & 0.000--0.550 \\
Protocol-semantic / app-layer & SqlInjection, XSS, DictionaryBruteForce, & \\
 & MITM-ArpSpoofing, DNS\_Spoofing, Recon-PortScan, & 0.000--0.025 \\
 & Uploading\_Attack, Mirai-\{greeth,greip,udpplain\} & \\
\bottomrule
\end{tabular}
\caption{Per-class correct-citation accuracy, grouped by observed pattern,
pooled across both models and both strategies.}
\label{tab:byclass}
\end{table}

\subsection{Cost}

Total real spend across the 5{,}053 scored trials: \$22.13. Claude
full-corpus (\$13.57) and RAG (\$7.85) together account for \$21.42
(96.8\%) of total spend; GPT-4o-mini full-corpus (\$0.46) and RAG (\$0.25)
together cost \$0.71. This matches the pre-registered cost estimate
(Section~\ref{sec:dataset}) of approximately \$21 for a complete run,
confirming the empirically-grounded per-call cost model carried over from
this project's earlier natural-language experiment generalizes to the
raw-feature task despite the different query content -- expected, since
input tokens are corpus-dominated (the fixed 15-entry corpus text is far
larger than either a short English sentence or a serialized 39-feature
row).

\section{Discussion}
\label{sec:disc}

The single largest finding is not the Claude-vs-GPT-4o-mini gap but the
per-class split (Section~\ref{sec:results}): a 39-feature flow summary
alone is apparently sufficient for near-perfect classification of
volumetric flood attacks, and close to useless for protocol-semantic and
application-layer attacks. This is a plausible, mechanistically sensible
result rather than a surprising one: a SYN flood's signature (extreme SYN
flag counts, extreme packet rate, near-zero completed handshakes) is
directly visible in the summary statistics this dataset extracts; a SQL
injection or XSS payload's signature is in the packet's application-layer
content, which these 39 features do not capture at all. The near-zero
accuracy on Mirai infection vectors (which in the real corpus are
grounded in device-compromise CVEs about default credentials and RCE, not
about the flood traffic Mirai devices later generate) makes the same point
from another angle: the corpus entries describe the \emph{infection}
mechanism, but the sampled flow features describe \emph{traffic}, a
category mismatch no amount of prompting fixes.

This reframes the Claude-vs-GPT-4o-mini comparison (Section~\ref{sec:results}):
Claude's significant advantage (52\% vs.\ 37\%, both strategies) is real
and consistent, but neither model is being tested primarily on "citation
skill" in the abstract -- both are being tested on whether they can
recognize volumetric-attack signatures in raw numeric data, since that is
where almost all of the correct-citation mass in this corpus actually
lives. A production pipeline built on this exact corpus and exact raw-
feature representation should expect strong performance on flood-type
detection and should not rely on either model for protocol-semantic attack
classification from flow features alone -- that would require either a
richer feature set (packet payload, not just header/timing statistics) or
a different grounding corpus entirely.

The near-total abstention failure (0.3--2.1\% correct-abstain, both
models, both strategies) is the second major finding and is worse than
this project has measured before. In the earlier natural-language
experiment, abstention was measured on only 5 held-out queries per cell and
was too noisy to interpret; here, with 600 abstention-expected trials per
cell, the low rate is precise, not noisy. Both models, shown a raw feature
vector from e.g.\ \texttt{Benign\_Final} (normal traffic) or
\texttt{VulnerabilityScan} (a class outside the 15-entry corpus's
coverage), overwhelmingly forced a citation onto the nearest available
entry rather than declining. This is the specific failure mode
citation-grounding is meant to prevent, and on this raw-feature task it is
not being prevented in practice -- a materially more important limitation
than the correct-citation numbers alone would suggest, since it means the
pipeline's citations cannot currently be trusted as evidence that a real
match exists, only that the model produced \emph{some} answer.

\section{Limitations}
\label{sec:limitations}

\begin{itemize}
    \item \textbf{387 Claude trials are missing, non-randomly.} All 387
    failures share one cause (account usage-limit reached,
    \texttt{2026-09-01} reset) and are concentrated at the end of the
    manifest's alphabetical class-processing order: \texttt{SqlInjection},
    \texttt{Uploading\_Attack}, \texttt{VulnerabilityScan}, and
    \texttt{XSS} have \emph{zero} successful Claude trials in this report
    (both strategies), and \texttt{Recon-PortScan} has only 6--7 of 40. This
    means the Claude correct-citation numbers in Table~\ref{tab:results-main}
    are not an unbiased estimate across all 19 mapped classes -- they
    entirely exclude two of the protocol-semantic classes already shown to
    score near 0\% for GPT-4o-mini, which, if included, would very likely
    lower Claude's aggregate correct-citation rate further, not raise it.
    This is the single most important caveat on this report's headline
    numbers and is why this paper is explicitly scoped as a partial-dataset
    report rather than a completed one.
    \item \textbf{Corpus covers 19 of 34 real classes.} The 15-entry corpus
    is unchanged from prior work and was not expanded for this raw-feature
    task; the 15 unmapped classes are scored only on abstention, not on
    whether a more complete corpus would have let the model cite them
    correctly.
    \item \textbf{TF-IDF retrieval only.} The RAG condition reuses the
    existing TF-IDF, $k{=}5$ retriever, not a neural embedding model; recall
    at this retriever's ceiling was already characterized in prior work and
    is not re-measured here.
    \item \textbf{Single sample, single seed.} $N{=}40$/class was drawn
    once with a fixed seed; no repeat-sampling variance estimate exists for
    the per-class accuracy figures in Table~\ref{tab:byclass}.
    \item \textbf{Feature set is header/timing statistics only.} The 39
    features extracted by this dataset's pipeline contain no
    application-layer payload content; the near-zero accuracy on
    protocol-semantic classes (Section~\ref{sec:disc}) may be a ceiling
    imposed by the feature extraction step, not by either model's
    reasoning ability, and this experiment cannot distinguish those two
    explanations.
    \item \textbf{Claude did not run at temperature 0.} Carried over from
    prior work: the Anthropic API rejects an explicit \texttt{temperature}
    parameter for \texttt{claude-sonnet-5}; GPT-4o-mini ran at true
    temperature 0. This asymmetry is a plausible, unmeasured confound on
    the model comparison.
    \item \textbf{Wrong-explanation rate not measured.} As in prior work,
    whether a correctly-cited entry's justification text is actually
    consistent with that entry's real content is not checked; a model could
    cite the right id for a fabricated reason and this would not be
    caught.
    \item \textbf{Single run, no repeated trials.} Each sample was sent
    once per (model, strategy) cell; no response-variability check was
    performed.
\end{itemize}

\section{Future Work}

In priority order: (1) re-run \texttt{csv\_flow\_experiment.py} once
Claude API access resets (2026-09-01) to backfill the 387 missing trials --
the harness is resumable and will retry only these, at no cost to the
already-collected data; (2) given the per-class split
(Section~\ref{sec:disc}), test whether a richer feature set (e.g.\
including payload-derived features, where available) closes the gap on
protocol-semantic classes, or whether that gap is fundamental to
header/timing-only features; (3) expand the grounding corpus to cover the
remaining 15 classes with real CAPEC/CVE entries, to distinguish "no
entry to cite" from "entry exists but the model didn't find it" for those
classes; (4) investigate the near-total abstention failure directly, e.g.\
by testing whether an explicit "if uncertain, abstain" instruction change
improves the 0.3--2.1\% rate; (5) repeat-sample multiple $N{=}40$ draws to
put a variance estimate on the per-class accuracy figures; (6) extend this
raw-feature protocol to the K3s CIS Self-Assessment Guide corpus under
separate development in this project.

\section{Conclusion}

Removing the natural-language intermediary between raw evidence and the
classifying model reveals two things the earlier, paraphrase-based version
of this experiment could not: first, that this project's citation
architecture, on raw 39-feature flow data, currently works well only for
volumetric flood attacks and not at all for protocol-semantic or
application-layer attacks, a mechanistic limit of the feature set rather
than an arbitrary model weakness; second, that correct abstention -- the
specific behavior citation-grounding exists to encourage -- is rare on this
harder task (0.3--2.1\%), a materially worse result than this project has
measured before and the clearest priority for future work. Claude
significantly outperforms GPT-4o-mini on raw-feature correct-citation rate
at both context strategies, but this report is explicit that its Claude
numbers exclude four real classes entirely due to a disclosed, non-random
data-collection gap (Section~\ref{sec:limitations}), and should be revised
once that gap is closed rather than treated as final.

\section*{Reproducibility Statement}

Dataset: CIC-IoT-2023 CSV release, 34 classes, 309 files, 8.4~GB,
46{,}776{,}697 rows, obtained and verified locally
(Section~\ref{sec:dataset}). Sampling: \texttt{build\_flow\_sample.py
--n-per-class 40 --seed 42}, producing \texttt{flow\_sample\_manifest.csv}
(1{,}360 rows, frozen). Corpus and ground-truth mapping:
\texttt{csv\_flow\_common.py}, reusing \texttt{real\_corpus.py}'s existing
\texttt{attack\_families} field unchanged from prior work
(Section~\ref{sec:corpus}). Experiment: \texttt{csv\_flow\_experiment.py},
resumable/checkpointed execution against live \texttt{claude-sonnet-5} and
\texttt{gpt-4o-mini} API calls, August 2026, per-trial output logged to
\texttt{csv\_flow\_results.csv}. Analysis:
\texttt{csv\_flow\_analysis.py}, producing \texttt{csv\_flow\_summary.csv}
and \texttt{csv\_flow\_by\_class.csv} and the McNemar's-test output quoted
in Section~\ref{sec:results}; this script explicitly excludes both
API-error rows and any harness self-test (\texttt{--mock}) rows sharing the
same output file from all real metrics, after both were found, during this
paper's preparation, to silently distort results if left in (mock rows
inflating apparent sample counts; error rows otherwise counted as incorrect
citations). This report covers the 5{,}053 real trials that completed
successfully as of the date in the title; 387 Claude trials remain pending
an account usage-limit reset and are explicitly out of scope
(Section~\ref{sec:limitations}), not silently included or extrapolated.

\begin{thebibliography}{9}

\bibitem[Bhaskaran et al.(2026)]{clearsight2026}
Bhaskaran, A., Ponnavaikko, P., and Shiva, G.\ S.
\textit{ClearSight: Citation-Grounded, Hallucination-Reduced IoT/OT Threat
Detection}. 2026.

\bibitem[Neto et al.(2023)]{neto2023ciciot2023}
Neto, E.~C.~P., Dadkhah, S., Ferreira, R., Zohourian, A., Lu, R., and
Ghorbani, A.~A.
CICIoT2023: A real-time dataset and benchmark for large-scale attacks in
IoT environment.
\textit{Sensors}, 23(13), 5941, 2023.

\bibitem[Majed et al.(2025)]{netflowviz}
Majed, L., Layeghy, S., Pandey, Y., Kulatilleke, G., and Portmann, M.
Multimodal LLMs for zero-shot intrusion detection using NetFlow
visualisations.
In \textit{Proc.\ IEEE 50th Conference on Local Computer Networks (LCN)},
2025.

\bibitem[Chandola et al.(2009)]{chandola2009anomaly}
Chandola, V., Banerjee, A., and Kumar, V.
Anomaly detection: A survey.
\textit{ACM Computing Surveys}, 41(3), 1--58, 2009.

\bibitem[Robertson and Zaragoza(2009)]{robertson2009bm25}
Robertson, S. and Zaragoza, H.
The Probabilistic Relevance Framework: BM25 and Beyond.
\textit{Foundations and Trends in Information Retrieval}, 3(4), 2009.

\bibitem[Deerwester et al.(1990)]{deerwester1990lsa}
Deerwester, S., Dumais, S.\ T., Furnas, G.\ W., Landauer, T.\ K., and
Harshman, R.
Indexing by Latent Semantic Analysis.
\textit{Journal of the American Society for Information Science}, 41(6),
1990.

\bibitem[MITRE(2026)]{capec}
MITRE Corporation. Common Attack Pattern Enumeration and Classification
(CAPEC). \url{https://capec.mitre.org}, accessed 2026.

\bibitem[NIST(2026)]{nvd}
National Institute of Standards and Technology. National Vulnerability
Database (NVD). \url{https://nvd.nist.gov}, accessed 2026.

\end{thebibliography}

\end{document}
